\documentclass[12pt]{article}
\usepackage[utf8]{inputenc}
\usepackage[T1]{fontenc}
\usepackage[margin=1in]{geometry}
\usepackage{setspace}
\usepackage{xcolor}
\usepackage[colorlinks=true, urlcolor=blue!60!black, linkcolor=blue!60!black]{hyperref}
\setlength{\parindent}{2em}
\setlength{\parskip}{1.25em}
\renewcommand{\baselinestretch}{1.0}
\begin{document}
\begin{center}
{\LARGE \textit{DeShaney v. Winnebago County Department of Social Services}, 489 U.S. 189 (1989)\par}
\vspace{0.8em}
{\large Case Brief --- Four-Page Edition\par}
\end{center}
\vspace{1em}

\section{Memory Jogger}

A county social-services agency repeatedly documented signs that four-year-old Joshua DeShaney was being abused by his father, yet the Court held 6--3 that the Due Process Clause generally imposes no affirmative duty to protect against private violence and that a brief earlier custody did not create a continuing duty; the State neither created the danger nor made Joshua more vulnerable (\href{https://supreme.justia.com/cases/federal/us/489/189/}{489 U.S. 189}).

\section{Detailed Case Facts}

After his parents' 1980 divorce, Joshua DeShaney lived with his father, Randy, in Winnebago County, Wisconsin. In 1982 Randy's second wife told police he had ``hit the boy, causing marks, and [was] a prime case for child abuse''; the Department of Social Services (DSS) accepted Randy's denial. In January 1983 Joshua was hospitalized with bruises; DSS obtained a juvenile court order placing him in the hospital's temporary custody, but three days later a ``Child Protection Team'' found the evidence insufficient and returned him under a voluntary agreement.

Over the next year, emergency room staff twice reported suspicious injuries, and a caseworker on monthly visits saw recurring head injuries, noted Randy's noncompliance, and recorded her suspicions, yet DSS did nothing, even when she was twice told Joshua was too ill to see her. In March 1984 Randy beat Joshua into a coma, leaving him permanently brain-damaged and institutionalized for life; Randy was convicted of child abuse.

\section{Procedural History}

Joshua and his mother sued the county, DSS, and its employees under 42 U.S.C. \S~1983 in the Eastern District of Wisconsin, alleging that their failure to intervene violated Joshua's substantive due process liberty interest. The district court granted summary judgment for respondents, the Seventh Circuit affirmed, and the Supreme Court, resolving a circuit split, affirmed on February 22, 1989.

\section{Judicial Votes}

Affirmed, \textbf{6--3}.
\begin{itemize}
\item \textbf{Majority:} Chief Justice Rehnquist, joined by Justices White, Stevens, O'Connor, Scalia, and Kennedy.
\item \textbf{Dissents:} Justice Brennan, joined by Justices Marshall and Blackmun; Justice Blackmun, separately.
\end{itemize}

\section{Holding}

The Due Process Clause is a limitation on the State's power to act, not a guarantee of minimal safety and security. It confers no affirmative right to protection against private violence unless the State has restrained the individual's liberty through custody, creating a ``special relationship,'' or has itself created or worsened the danger. DSS's knowledge and monitoring did not create such a duty, and returning Joshua to his father placed him in no worse position.

\section{Analysis of Opinions}

\textbf{Majority (Rehnquist): negative liberties.} The text forbids the State to ``deprive'' persons of life, liberty, or property. Petitioners claimed a substantive entitlement, having waived any procedural-entitlement claim, and neither text nor history supports a duty to prevent private harm; the scope of protective services is left to democratic politics. Under \textit{Harris v. McRae}, a State with no duty to provide services cannot be liable for injuries they might have averted, absent selective denial violating Equal Protection.

\textbf{The special-relationship exception.} \textit{Estelle v. Gamble} and \textit{Youngberg v. Romeo} hold only that when the State holds a person in custody against his will, rendering him unable to care for himself, its restraint triggers a duty to meet basic needs. Joshua was harmed in his father's custody, not the State's; the Court reserved whether state foster care is analogous. Returning him ``placed him in no worse position than that in which he would have been had it not acted at all.''

\textbf{Tort law and the social-worker dilemma.} The Amendment ``does not transform every tort committed by a state actor into a constitutional violation''; Wisconsin may impose tort liability, and removing Joshua prematurely would have invited due process claims by the parent.

\textbf{Brennan (dissenting).} The majority's action--inaction baseline ignored Wisconsin's affirmative choices: its child-welfare system funneled all abuse reports to DSS, displacing other rescuers, and the caseworker later said, ``I just knew the phone would ring some day and Joshua would be dead.'' By monopolizing protection and then failing to act, DSS confined Joshua in a violent home.

\textbf{Blackmun (dissenting).} Condemning ``sterile formalism'' and comparing it to antebellum judges enforcing the Fugitive Slave Acts, Blackmun urged a ``sympathetic'' reading and lamented, ``Poor Joshua!''

\textbf{Later developments.} \textit{Collins v. City of Harker Heights} (1992) and \textit{Town of Castle Rock v. Gonzales} (2005) reinforced the rule, while the ``played no part'' sentence spawned the circuits' state-created-danger doctrine, never formally endorsed by the Supreme Court.

\section{Examples: Future Applications}

\begin{enumerate}
\item \textbf{Same-Side (State Prevails): Slow 911 Response.} Police log three calls about a protective-order violation as low priority, and the spouse is assaulted. Under \textit{DeShaney} and \textit{Castle Rock}, no custody, no created danger, no claim.
\item \textbf{Same-Side (State Prevails): Broken Patrol Promise.} A woman threatened by a former boyfriend is injured after an officer fails to deliver promised hourly patrols. A promise is among the State's ``expressions of intent to help him,'' not a restraint on liberty.
\item \textbf{Opposite-Side (Potential Claim): Child Stranded on a Highway.} A trooper orders a seven-year-old out of a car onto an unlit interstate at 1:00 a.m. The State created the peril (cf.\ \textit{White v. Rochford}).
\item \textbf{Opposite-Side (Duty): Prisoner and Cellmate.} Custody triggers a duty of reasonable safety under \textit{Estelle}, \textit{Youngberg}, and \textit{Farmer v. Brennan}.
\item \textbf{Fence-Sitter: State Foster Care.} A child placed by the State in a licensed foster home is abused after ignored reports; footnote 9 left this open.
\end{enumerate}

\section{Critique}

\textbf{Action versus inaction formalism.} Critics echo Brennan that ``inaction'' is a baseline choice: Wisconsin affirmatively enacted custody laws giving Randy exclusive control over Joshua and a reporting system that funneled every rescue effort to DSS.

\textbf{Reconstruction history.} \textit{DeShaney} stands in tension with the Fourteenth Amendment's origins, enacted largely because Southern states refused to protect Black citizens from Ku Klux Klan violence, though Rehnquist addressed selective non-protection only under the Equal Protection Clause.

\section{Key Quotations}

\begin{quote}
``While the State may have been aware of the dangers that Joshua faced in the free world, it played no part in their creation, nor did it do anything to render him any more vulnerable to them.''
\end{quote}
\noindent --- Chief Justice Rehnquist, for the Court, 489 U.S. at 201.

\begin{quote}
``The Clause is phrased as a limitation on the State's power to act, not as a guarantee of certain minimal levels of safety and security. It forbids the State itself to deprive individuals of life, liberty, or property without `due process of law,' but its language cannot fairly be extended to impose an affirmative obligation on the State to ensure that those interests do not come to harm through other means. . . . Its purpose was to protect the people from the State, not to ensure that the State protected them from each other.''
\end{quote}
\noindent --- Chief Justice Rehnquist, for the Court, 489 U.S. at 195--196.

\begin{quote}
``The affirmative duty to protect arises not from the State's knowledge of the individual's predicament or from its expressions of intent to help him, but from the limitation which it has imposed on his freedom to act on his own behalf.''
\end{quote}
\noindent --- Chief Justice Rehnquist, for the Court, 489 U.S. at 200.

\end{document}
